% Applied Software Development Lab -- Lab Handout 9: More Control!
% Copyright (c) Michael Lehn. Licensed under CC BY-SA 4.0.
%
% Region coordinates (bp from the top of the page) follow the page layout of
% the original German handout "Laborblatt 9".
\documentclass{aslab}
\sethandout{9}{More Control!}
\graphicspath{{../figures/handout09/}{../figures/shared/}}

% heading color of the "motivation" box on page 7
\definecolor{motivmagenta}{HTML}{970E53}
% code lines keep the region pitch even with tall or deep glyphs
\makeatletter
\lst@AddToHook{Init}{\lineskiplimit=-\maxdimen\relax}
\makeatother
% "•" lists with the text indented by 9bp
\newlist{dotlist}{itemize}{1}
\setlist[dotlist]{label={\textbullet},leftmargin=9bp,labelsep=*,align=left,
  labelindent=0pt}
% inline code with leading spaces (four spaces = 2.4em in the mono font)
\newcommand*\indentcode[2]{\codearg{\hspace*{#1em}#2}}

\begin{document}
%% ===================================================================== 1 ==
\maketitleblock

\begin{region}[pitch=13]{124.0}
\heading{Learning objectives}
\end{region}

\begin{region}[pitch=13]{138.5}
{\bfseries\color{aslabblue}
\begin{dashlist}
\item Connecting hardware (push button \& potentiometer) to a game via the
  serial interface
\item Understanding and reusing existing microcontroller code
\item Establishing the connection between game and controller
\item How to write a ``proper'' module
\end{dashlist}
}
\end{region}

\begin{region}{218.8}
\heading{Looking back and today's goal}
\end{region}

\begin{region}[pitch=14.1]{246.4}
So far, you have worked on the Breakout game, which was controlled with the
keyboard. Today we extend it -- with a self-built \textbf{game controller on
the breadboard}!
\end{region}

\begin{region}[pitch=14.05]{287.9}
In \textbf{lab handouts 4 and 5} you programmed a microcontroller that
\begin{dashlist}
\item detects the \textbf{push-button state} (pressed / released)
\item measures the \textbf{analog value} of the potentiometer (a
  \textbf{10-bit value}, i.e., an integer in the range \textbf{from 0 to
  1023})
\end{dashlist}
This information is sent via the serial interface -- for example:
\end{region}

\begin{region}[x=68.7,right=0,pitch=14.25]{357.2}
\begin{dashlist}
\item \textbf{Push button pressed}: a line with \code{A}
\item \textbf{Push button released}: a line with \code{a}
\item \textbf{Voltage change at the potentiometer}: a line such as
  \code{X 528}
\end{dashlist}
\end{region}

\begin{region}[pitch=13.65]{400.1}
You can now build this data directly into the game:
\begin{dashlist}
\item The \textbf{push button} fires rockets
\item The \textbf{potentiometer} controls the \textbf{position of the paddle}
\end{dashlist}
\end{region}

\begin{region}{470.2}
\heading{Preparation}
\end{region}

\begin{region}[pitch=13.7]{497.7}
As usual, we work with two terminals:
\begin{dashlist}
\item \textbf{left}: editing source code
\item \textbf{right}: compiling code, flashing, opening the serial monitor,
  testing the game
\end{dashlist}
\end{region}

\begin{region}{566.7}
\subheading{Cloning the repository}
\end{region}

\begin{region}{593.2}
In the left or right terminal:
\end{region}

\begin{region}{619.8}
\code{git clone https://gitlab.uni-ulm.de/USERNAME/asp-ulm}
\end{region}

\begin{region}{646.2}
Then, in both terminals, change into the directory with the microcontroller
code:
\end{region}

\begin{region}{672.8}
\code{cd asp-ulm/controller}
\end{region}

\begin{region}[pitch=13]{698.9}
In the next step, we take a closer look at the existing code once more:
\begin{dashlist}
\item How does the Makefile work?
\item How is the code compiled and transferred to the microcontroller?
\item How do you open the serial monitor to see the data being sent?
\end{dashlist}
\end{region}

\nextpage
%% ===================================================================== 2 ==
\begin{region}{69.4}
\heading{Step 1: Microcontroller code and breadboard tinkering}
\end{region}

\begin{region}[pitch=13.1]{96.9}
In the \textbf{left terminal}, open the file \code{controller.ino} containing
the C++ code for the microcontroller. Go through the code once more at your
own pace and try to understand exactly what happens. If anything is
unclear, you can of course
\begin{dashlist}
\item take another look at lab handout 4,
\item ask us,
\item or help each other out.
\end{dashlist}
\end{region}

\begin{region}{180.4}
By the way, we like to ask silly questions ourselves now and then -- that is
part of it \emoji{1F609}
\end{region}

\begin{region}[pitch=14.27]{210.3}
Before you run the microcontroller code, you should rebuild \textbf{the
familiar circuit} with \textbf{push button and potentiometer on the
breadboard} and connect it correctly to the microcontroller. Here is a
picture of the complete setup (the push button is connected to the digital
input \textbf{pin~2}, the potentiometer to the analog input \textbf{A0}):
\end{region}

\placegraphic{96.0}{257.1}{403.3}{breadboard-setup.png}

\begin{region}[pitch=14.0]{470.4}
The current directory \code{controller} also contains a file named
\code{Makefile}. If you like, you can briefly display its (somewhat cryptic)
content in the \textbf{right terminal} with \code{cat Makefile}. For today,
however, it is enough to know \textbf{what this Makefile is good for}. It
contains rules for automatically invoking the right tools -- namely for:
\end{region}

\begin{region}[pitch=13.4]{525.3}
\begin{dashlist}
\item the \textbf{translation} of the C++ code (\code{controller.ino}) into
  \textbf{machine code} for the microcontroller
\item flashing the microcontroller with avrdude
\item opening a serial monitor to see the output of the controller
\end{dashlist}
\end{region}

\begin{region}{592.6}
The program \code{make} does all of this for you -- depending on the target
you invoke:
\end{region}

\begin{region}[x=60.7,right=0,pitch=22.73]{611.1}
\newcommand\cmdcol[1]{\makebox[122.1bp][l]{#1}}%
\cmdcol{\small\bfseries Command}{\small\bfseries Meaning}\par
\cmdcol{\codearg{make}}{\small Translates the code into machine code (hex
  file)}\par
\cmdcol{\codearg{make upload}}{\small Flashes the microcontroller (compiles
  beforehand if necessary)}\par
\cmdcol{\codearg{make monitor}}{\small Opens the serial monitor, which shows
  what the microcontroller sends}
\end{region}

\begin{region}[pitch=15.15]{704.3}
\emoji{26A0} Important: For \code{make upload} or \code{make monitor} to
work, \textbf{the microcontroller must be connected} and \textbf{the USB
port must be made visible under WSL} -- as usual with: \code{usb_check}
\end{region}

\begin{region}[pitch=13.3]{762.0}
\textbf{Task:} Once everything is set up and the microcontroller is
connected, compile the program, flash the microcontroller, and use the
serial monitor to check whether everything works as expected.
\end{region}

\nextpage
%% ===================================================================== 3 ==
\begin{region}{69.4}
\heading{Step 2: Reading the serial interface with Python
  (\codearg{\mdseries serial-reader.py})}
\end{region}

\begin{region}[pitch=13.2]{97.0}
By the way, we also wrote our own serial monitor in Python (see lab
handout~4). Let's take another look at it now -- this time with fresh eyes!
\end{region}

\begin{region}[pitch=13.6]{137.5}
\subheading{Preparation}
In both terminals, change into the neighboring directory \code{python}:
\end{region}

\begin{region}{177.8}
\code{cd ../python}
\end{region}

\begin{region}[pitch=14.3]{205.1}
In the \textbf{left terminal}, open the file \code{serial-reader.py}.\\
In the \textbf{right terminal}, you can start the monitor with:
\end{region}

\begin{region}{246.1}
\code{python serial-reader.py}
\end{region}

\begin{region}{273.5}
\subheading{Task}
\end{region}

\begin{region}[pitch=13.3]{301.0}
When we worked on \textbf{lab handout 4}, we were not yet very familiar with
object-oriented programming. That has changed now -- and that is exactly
why a second look is worthwhile!
\end{region}

\begin{region}[pitch=13.2]{340.9}
Go through the code \code{serial-reader.py} step by step -- and work on the
following tasks:
\end{region}

\begin{region}{381.4}
\begin{tasks}
\item \textbf{Object-oriented interface to the serial connection}
\end{tasks}
\end{region}

\begin{region}[x=74.7,right=0,pitch=13.5]{399.4}
\code{ser = serial.Serial("/dev/ttyUSB0", 9600)}\\
Here, an object of the class \code{Serial} from the module \code{serial} is
created. This object (which we call \code{ser}) manages:
\end{region}

\begin{region}[x=74.7,right=0,pitch=13.15]{441.1}
\begin{dotlist}
\item the serial interface,
\item an (internal) buffer for received bytes,
\item and provides methods for working with it.
\end{dotlist}
\end{region}

\begin{region}[x=74.7,right=0,pitch=14.3]{479.2}
Important methods:
\begin{dotlist}
\item \code{in_waiting}: indicates \textbf{how many bytes are currently in
  the buffer}
\item \code{read(n)}: reads \textbf{n bytes from the buffer} (FIFO -- first
  in, first out)
\end{dotlist}
In our code, we always read \textbf{a single byte} with \code{ser.read(1)}.
\end{region}

\begin{region}{549.3}
\begin{tasks}[start=2]
\item \textbf{Small code change}
\end{tasks}
\end{region}

\begin{region}[x=74.7,right=0,pitch=13.35]{567.0}
In the code you will find this line:\\
\code{ch = bytes.decode(ch)}\\
Change it to:\\
\code{ch = ch.decode()}\\
Then start the program again and check whether everything still works as
before. Discuss among yourselves, or ask us, why the new line is equivalent
to the old one.
\end{region}

\begin{region}[x=74.7,right=0]{660.8}
\itshape Note: \codearg{\upshape ser.read(1)} returns an object of type
\codearg{\upshape bytes}.
\end{region}

\begin{region}[pitch=13.2]{687.3}
When you are done, make a commit so that the small change to the code is not
lost:
\end{region}

\begin{region}[pitch=13.2]{727.1}
\code{git commit -a -m "call decode as a method"}\\
\code{git push}
\end{region}

\nextpage
%% ===================================================================== 4 ==
\begin{region}{69.4}
\heading{Step 3: Python basics (splitting strings -- converting numbers)}
\end{region}

\begin{region}[pitch=13.2]{97.0}
Before we build the serial data into the game, we still need a few practical
basics in Python.
\end{region}

\begin{region}[pitch=13.4]{137.5}
\textbf{The problem:} The serial interface delivers strings such as
\code{"X 512"}. It would be more practical if we could turn them directly
into a list such as \code{["X", 512]} -- that is: split the text at the
space and convert the second part into a number.
\end{region}

\begin{region}{191.6}
\subheading{A short Python session}
\end{region}

\begin{region}{218.4}
In the right terminal, start the Python shell with: \code{python}. Then off
you go:
\end{region}

\begin{region}[pitch=13.4]{244.8}
\begin{tasks}
\item Enter:\\
  \tcode{>>> }\code{s = "X 512"}
\end{tasks}
\end{region}

\begin{region}[pitch=13.4]{271.2}
\begin{tasks}[start=2]
\item Now split the string:\\
  \tcode{>>> }\code{l = s.split()}
\end{tasks}
\end{region}

\begin{region}[pitch=13.42]{297.5}
\begin{tasks}[start=3]
\item Check the content:\\
  \tcode{>>> }\code{s}\\
  \tcode{>>> }\code{l}\\
  $\rightarrow$ \code{l} should contain \code{["X", "512"]}.\\
  Also try:\\
  \tcode{>>> }\code{s = "X    512"}\\
  \tcode{>>> }\code{l = s.split()}\\
  Several spaces make no difference -- \code{split()} ignores superfluous
  spaces.
\end{tasks}
\end{region}

\begin{region}[pitch=13.35]{418.6}
\begin{tasks}[start=4]
\item Access individual list elements:\\
  \tcode{>>> }\code{l[0]}\\
  \tcode{>>> }\code{l[1]}
\end{tasks}
\end{region}

\begin{region}[pitch=13.45]{458.2}
\begin{tasks}[start=5]
\item What happens here?\\
  \tcode{>>> }\code{l[1] + 1}\\
  This produces an error -- because \code{l[1]} is a string, not a number.
\end{tasks}
\end{region}

\begin{region}[pitch=13.3]{498.4}
\begin{tasks}[start=6]
\item Convert the string into a number and try again:\\
  \tcode{>>> }\code{l[1] = int(l[1])}\\
  \tcode{>>> }\code{l[1] + 1}
\end{tasks}
\end{region}

\begin{region}[pitch=13.28]{538.0}
\begin{tasks}[start=7]
\item We could also process strings with more than one number -- for
  example:\\
  \tcode{>>> }\code{s = "X 523 23"}\\
  \tcode{>>> }\code{l = s.split()}\\
  \tcode{>>> }\code{for i in range(1, len(l)):}\\
  \tcode{... }\indentcode{2.4}{l[i] = int(l[i])}\\
  $\rightarrow$ Then display \tcode{l}.\\
  You will see: from the second element on, all entries have been converted
  into integers.
\end{tasks}
\end{region}

\begin{region}[x=74.7,right=0,pitch=13.2]{646.2}
\subheading{Why this is useful}
Normally, we only want to program what is necessary -- no generalizations in
advance. But here we actually do need this generalization -- not because of
cases with several numbers, but because of cases with no number at all:\\
\tcode{>>> }\code{s = "A"}\\
\tcode{>>> }\code{l = s.split()}\\
\tcode{>>> }\code{for i in range(1, len(l)):}\\
\tcode{... }\indentcode{2.4}{l[i] = int(l[i])}\\
This works as well -- because the loop is not executed at all when \code{l}
has only one element.
\end{region}

\nextpage
%% ===================================================================== 5 ==
\begin{region}{69.4}
\heading{Step 4: Our Python interface for the serial interface
  (\codearg{\mdseries controller.py})}
\end{region}

\begin{region}[pitch=14.3]{96.9}
In the \textbf{left terminal}, open the file \code{controller.py} (it was
created in lab handout 5)\\
and start it in the \textbf{right terminal} with: \code{python controller.py}
\end{region}

\begin{region}[pitch=14.3]{138.6}
\subheading{Is this object-oriented?}
Although a class Controller appears in the code, this is \textbf{not} an
object-oriented interface -- which is quite common for hardware-oriented
interfaces.
\end{region}

\begin{region}{185.8}
Instead, there are \textbf{two free functions}, which can be used like this:
\end{region}

\begin{region}[pitch=13.2]{204.9}
\tcode{connect("/dev/ttyUSB0")}\\
\tcode{line = readline()}
\end{region}

\begin{region}[pitch=13.63]{236.4}
\begin{dashlist}
\item \code{connect(port)} must be called first to establish the connection.
\item \code{readline()} returns either:
  \begin{dashlist}
  \item \code{None} (if no complete line is available yet), or
  \item the \textbf{received string} (if a line was terminated with
    \code{\n}).
  \end{dashlist}
\end{dashlist}
\end{region}

\begin{region}[pitch=13.6]{305.7}
\subheading{What does the class actually do?}
The class \code{Controller} contains two static attributes:
\begin{dashlist}
\item \code{Controller.ser} -- the serial connection
\item \code{Controller.buffer} -- the current input buffer
\end{dashlist}
\end{region}

\begin{region}[pitch=13.23]{360.5}
They serve as \textbf{global variables}, but are \textbf{clearly
encapsulated} by the class -- which shows that they belong together
logically.\\
The two functions connect() and readline() access these attributes -- so
you could also move them directly into the class as static methods and then
call them like this:
\end{region}

\begin{region}[pitch=13.2]{419.1}
\tcode{Controller.connect(port)}\\
\tcode{Controller.readline()}
\end{region}

\begin{region}{451.7}
You \textbf{don't have to} do this -- but you \textbf{may}, if you feel like
trying it out.
\end{region}

\begin{region}{479.1}
\subheading{Task}
\end{region}

\begin{region}{497.9}
\begin{tasks}
\item Now write a function:
\end{tasks}
\end{region}

\begin{region}[x=74.7,right=0]{516.9}
\tcode{def readdata():}
\end{region}

\begin{region}[x=74.7,right=0,pitch=13.5]{535.9}
It should try to read a line with \code{readline()}. If that was successful
(i.e., you did not get \code{None}):
\end{region}

\begin{region}[x=74.7,right=0,pitch=13.6]{564.6}
\begin{dotlist}
\item split the line into a list as in \textbf{Step 3}, i.e., \textbf{convert
  all numbers from the second element on} with \code{int()},
\item and return the resulting list.
\end{dotlist}
\end{region}

\begin{region}[x=74.7,right=0,pitch=13.5]{603.8}
If no line is available (\code{readline()} returned \code{None}),
\code{readdata()} should also return the value \code{None}.
\end{region}

\begin{region}[x=74.7,right=0]{643.7}
To test it, in the main loop of the program simply replace
\end{region}

\begin{region}[x=74.7,right=0]{662.7}
\tcode{line = readline()}
\end{region}

\begin{region}[x=74.7,right=0]{681.4}
with
\end{region}

\begin{region}[x=74.7,right=0]{700.3}
\tcode{line = readdata()}
\end{region}

\begin{region}[x=74.7,right=0,pitch=13.4]{719.1}
Then restart the program and check whether everything works as expected --
that is: instead of \code{"A"} or \code{"a"} you get \code{['A']} or
\code{['a']}, respectively, and for \code{"X 523"} you get
\code{['X', 523]}.
\end{region}

\begin{region}[pitch=13.5]{745.7}
\begin{tasks}[start=2]
\item If everything works:\\
  \code{git commit -a -m "serial parser: readdata"}\\[-0.3bp]
  \code{git push}
\end{tasks}
\end{region}

\nextpage
%% ===================================================================== 6 ==
\begin{region}{69.4}
\heading{Step 5: Turning \codearg{\mdseries controller.py} into a module}
\end{region}

\begin{region}[pitch=13.1]{87.2}
So that we can later use our code from other Python programs (e.g., from
within the game), we turn controller.py into a proper module.
\end{region}

\begin{region}[pitch=13.2]{117.9}
To do this, we move the previous code from the main loop into a function
main() and add the following at the end of the file:
\end{region}

\begin{region}[pitch=13.2]{150.0}
\begin{codeblock}
if __name__ == "__main__":
    main()
\end{codeblock}
\end{region}

\begin{region}{182.6}
\subheading{What is the benefit?}
\end{region}

\begin{region}[pitch=14.4]{201.3}
\begin{dashlist}
\item If you \textbf{run} \code{controller.py} \textbf{directly} with
  \code{python controller.py}, the \code{main()} function is called as
  before and the program runs. So you can still \textbf{test} the module.
\end{dashlist}
\end{region}

\begin{region}[pitch=14.3]{243.4}
\begin{dashlist}
\item If you \textbf{import} the module, e.g., with \code{import controller},
  \textbf{nothing is executed automatically}. So you can build the
  functionality into other places selectively and use it there. That is
  exactly what we do in the next step!
\end{dashlist}
\end{region}

\begin{region}[pitch=18.8]{299.3}
\subheading{Task}
Quickly test whether everything still works as usual when you run python
controller.py. If so:
\end{region}

\begin{region}[pitch=13.2]{337.0}
\code{git commit -a -m "turned into module"}\\
\code{git push}
\end{region}

\begin{region}{380.6}
\heading{Step 6: Firing rockets with the push button \emoji{1F680}}
\end{region}

\begin{region}[pitch=14.2]{397.8}
Now we want to test how to build the module \code{controller.py} into our
game -- using this example: \textbf{The push button fires a rocket.}
\end{region}

\begin{region}[pitch=14.3]{427.3}
\begin{tasks}
\item \textbf{Importing the module}\\
  In the \textbf{left terminal}, open the file \code{breakout.py}. At the
  top of the code, import the module with
\end{tasks}
\end{region}

\begin{region}[x=74.7,right=0]{459.6}
\tcode{import controller}
\end{region}

\begin{region}[x=74.7,right=0,pitch=13.5]{477.2}
If you implemented \code{connect}, \code{readline}, and \code{readdata} as
static methods within the class Controller, use:
\code{from controller import Controller}; otherwise you would have to write
\code{controller.Controller.connect} instead of \code{Controller.connect}.
\end{region}

\begin{region}[pitch=13.3]{518.4}
\begin{tasks}[start=2]
\item \textbf{Establishing the connection to the serial interface}\\
  Insert the following line before the main loop:
\end{tasks}
\end{region}

\begin{region}[x=74.7,right=0]{550.6}
\tcode{controller.connect("/dev/ttyUSB0")}
\end{region}

\begin{region}[x=74.7,right=0,pitch=13.4]{569.4}
If you implemented connect as a static method within the class Controller,
use:\\
\code{Controller.connect("/dev/ttyUSB0")}
\end{region}

\begin{region}[pitch=13.3]{597.0}
\begin{tasks}[start=3]
\item \textbf{Processing serial data in the game loop}\\
  Insert the following code in the main loop:
\end{tasks}
\end{region}

\begin{region}[x=74.7,right=0,pitch=13.2]{628.1}
\begin{codeblock}
ser_data = controller.readdata()
if ser_data:
    if ser_data[0] == "A":
       rockets.append(Rocket(Paddle.x + Paddle.w // 2, Paddle.y, -10))
\end{codeblock}
\end{region}

\begin{region}[x=74.7,right=0,pitch=13.5]{685.3}
Here, too: if you have \code{readdata()} as a static method in the class,
use\\
\code{ser_data = Controller.readdata()}
\end{region}

\begin{region}[x=74.7,right=0,pitch=13.4]{717.6}
The code checks whether a new line has been received from the serial
interface. If its first component is \code{"A"} (push button pressed), a
rocket is created and fired.
\end{region}

\begin{region}[pitch=13.3]{753.0}
Test your game. If everything works:\\
\code{git commit -a -m "button fires"}\\[-0.1bp]
\code{git push}
\end{region}

\nextpage
%% ===================================================================== 7 ==
\begin{region}{72.9}
\heading{Step 7: Controlling the paddle with the potentiometer
  \emoji{1F579}}
\end{region}

\begin{region}{96.2}
For now, only one thing is missing: \textbf{controlling the paddle directly
with the potentiometer.}
\end{region}

\begin{region}[pitch=13.3]{115.2}
For this, we extend the class Paddle with a method \code{set()} that lets
you set the \code{x} position directly.
\end{region}

\begin{region}[pitch=13.3]{150.3}
\begin{tasks}
\item \textbf{Extending the code in \codearg{\mdseries breakout.py}}\\
  Change the code that queries the controller to:
\end{tasks}
\end{region}

\begin{region}[x=74.7,right=0,pitch=13.26]{182.5}
\begin{codeblock}
ser_data = controller.readdata()
if ser_data:
   if ser_data[0] == "A":
      rockets.append(Rocket(Paddle.x + Paddle.w // 2, Paddle.y, -10))
   elif ser_data[0] == "X":
       Paddle.set(ser_data[1])
\end{codeblock}
\end{region}

\begin{region}[pitch=13.6]{268.5}
\begin{tasks}[start=2]
\item \textbf{Writing the method \codearg{\mdseries set()} in paddle.py}\\
  Open the file \code{paddle.py} and add the method:
\end{tasks}
\end{region}

\begin{region}[x=74.7,right=0,pitch=13.2]{301.1}
\begin{codeblock}
def set(x):
    Paddle.x = x
    if Paddle.x < 0:
        Paddle.x = 0
    if Paddle.x + Paddle.w > Game.width:
        Paddle.x = Game.width - Paddle.w
\end{codeblock}
\end{region}

\begin{region}[x=74.7,right=0]{385.8}
This sets \code{Paddle.x} and limits it to the left and right if necessary.
\end{region}

\begin{region}[pitch=13.25]{400.0}
\begin{tasks}[start=3]
\item \textbf{Testing and fine-tuning}\\
  Start the game and test whether you can move the paddle with the
  potentiometer.\\
  If it feels too sluggish, try in breakout.py:
\end{tasks}
\end{region}

\begin{region}[x=74.7,right=0,pitch=18.8]{445.4}
\tcode{Paddle.set(ser_data[1] * 1.5)}\\
or\\
\tcode{Paddle.set(ser_data[1] * 2)}
\end{region}

\begin{region}{509.2}
Once everything works (and you are happy with the paddle speed):
\end{region}

\begin{region}[pitch=13.2]{535.7}
\code{git commit -a -m "paddle works too"}\\
\code{git push}
\end{region}

\begin{region}{594.8}
{\fontsize{13bp}{15bp}\selectfont\bfseries\color{motivmagenta}Motivation for
the next step: it's about cold hard cash \emoji{1F4B0}\par}
\end{region}

\begin{region}{624.9}
Now, as a test, unplug the USB cable from the microcontroller and restart the
game.
\end{region}

\begin{region}[pitch=15.8]{642.4}
\emoji{1F4A5} The game crashes because the call
\code{controller.connect("/dev/ttyUSB0")} fails.
\end{region}

\begin{region}[pitch=13.15]{658.2}
Imagine: you have programmed a great game that is already fun without a
controller -- but only with a controller does it become really cool. And
that is exactly what people are willing to spend a few hard-earned euros on.
\end{region}

\begin{region}{715.2}
\emoji{1F449} For this to work, your game must:
\end{region}

\begin{region}[pitch=14.2]{730.9}
\begin{dashlist}
\item \textbf{run smoothly even without a controller},
\item[{\color{aslabblue}\fontsize{13.2bp}{12bp}\selectfont -}] and
  ideally even allow \textbf{connecting the controller while the game is
  running}.
\end{dashlist}
\end{region}

\nextpage
%% ===================================================================== 8 ==
\begin{region}{69.4}
\heading{Step 8: Catching exceptions}
\end{region}

\begin{region}[pitch=13.1]{90.5}
We now make our code more robust by cleanly catching errors that occur while
establishing the connection.
\end{region}

\begin{region}{125.4}
\subheading{Task}
\end{region}

\begin{region}[pitch=13.3]{145.2}
\begin{tasks}
\item \textbf{Adapting} \code{connect()}\\
  In controller.py, change the function connect as follows:
\end{tasks}
\end{region}

\begin{region}[x=74.7,right=0,pitch=13.2]{177.4}
\begin{codeblock}
def connect(port):
    try:
        Controller.ser = serial.Serial(port, 9600)
        Controller.buffer = ""
        print(f"connected to {port}")
    except:
        print(f"could not connected to {port}")
\end{codeblock}
\end{region}

\begin{region}[x=74.7,right=0,pitch=15.5]{276.0}
\subheading{What happens here?}
Here you see the basic pattern:
\end{region}

\begin{region}[x=74.7,right=0,pitch=13.2]{310.4}
\begin{codeblock}
try:
    ...
except:
    ...
\end{codeblock}
\end{region}

\begin{region}[x=74.7,right=0,pitch=13.5]{368.7}
This means: the program tries to execute the code in the \code{try} block.
If an error occurs, the program jumps to the \code{except} block. If no
error occurs, the \code{except} block is simply skipped.
\end{region}

\begin{region}[x=74.7,right=0,pitch=13.28]{414.5}
In our case, this specifically means:
\begin{dotlist}
\item If the serial interface can be opened, everything is fine --
  \code{Controller.ser} contains a valid object.
\item If it fails, there is just a message, and \code{Controller.ser} stays
  \code{None} -- the program does not crash (at this point).
\end{dotlist}
\end{region}

\begin{region}[pitch=13.8]{495.1}
\begin{tasks}[start=2]
\item \textbf{Testing and solving the next problem}\\
  Start the game \textbf{without a connected microcontroller}. You will see:
  it still crashes -- but only later, namely when \code{readline()} is
  called. Why? Because this function assumes that \code{Controller.ser} is
  always a valid object.
\end{tasks}
\end{region}

\begin{region}[x=74.7,right=0,pitch=13.5]{558.8}
This is easy to fix. We only need to safeguard \code{readline()}. At the
beginning of the function \code{readline()}, check whether
\code{Controller.ser} is a valid object:
\end{region}

\begin{region}[x=74.7,right=0,pitch=13.2]{591.3}
\begin{codeblock}
def readline():
    if Controller.ser is None:
        return None
    # rest as before
\end{codeblock}
\end{region}

\begin{region}[pitch=13.5]{650.2}
\begin{tasks}[start=3]
\item \textbf{Testing and solving the next problem}\\
  Please start the game once more, again without connecting the
  microcontroller. You will notice that it crashes again, but this time at a
  different place, namely in the function \code{readdata()}.\\
  This happens because we previously added a check to \code{readline()} that
  tests whether \code{Controller.ser} is a valid object at all. If it is
  not, \code{readline()} returns the value \code{None}. And that is exactly
  the problem: in the function \code{readdata()}, \code{s.split()} is called
  immediately without any further check. As long as \code{s} is a string,
  this works fine, but if \code{s} has the value \code{None}, the program
  crashes, because a \code{None} object has no \code{split()} method.
\end{tasks}
\end{region}

\nextpage
%% ===================================================================== 9 ==
\begin{region}[x=74.7,right=0,pitch=13.3]{81.9}
To prevent this, we need to adapt \code{readdata()} so that it first checks
whether a value is actually present. The function could then look like
this, for example:
\end{region}

\begin{region}[x=74.7,right=0,pitch=13.2]{121.8}
\begin{codeblock}
def readdata():
    s = readline()
    if s is not None:
        s = s.split()
        for i in range(1, len(s)):
            s[i] = int(s[i])
        return s
    return None
\end{codeblock}
\end{region}

\begin{region}[x=74.7,right=0]{240.3}
Now the game should also run without a connected microcontroller.
\end{region}

\begin{region}[pitch=13.4]{267.6}
\begin{tasks}[start=4]
\item \textbf{Test: connecting the microcontroller while the game is
  running}\\
  As an exception, the \textbf{left terminal} is now \textbf{not used for
  editing}, but for running \code{usb_check} -- our familiar workaround that
  lets WSL see the USB port. We want to test whether the game also works
  when the microcontroller is connected. We distinguish two cases:
\end{tasks}
\end{region}

\begin{region}[x=74.7,right=0,pitch=13.2]{335.0}
  \textbf{(a)} The microcontroller is already connected before the game
  starts, or\\
  \textbf{(b)} it is connected while the game is already running.
\end{region}

\begin{region}[x=74.7,right=0,pitch=13.4]{375.5}
\textbf{First, test case (a):} Connect the microcontroller, run
\code{usb_check} in the left terminal, and then start the game in the right
terminal with \code{python breakout.py}. You should now be able to control
the game with the breadboard.
\end{region}

\begin{region}[x=74.7,right=0,pitch=13.5]{429.6}
\textbf{Next, case (b):} Unplug the USB cable from the microcontroller,
restart the game in the right terminal with \code{python breakout.py}, then
reconnect the microcontroller and run \code{usb_check} in the left
terminal.\\
You will notice: \textbf{unfortunately, the breadboard no longer works
now}. The reason is simple: our program tries only once, at the beginning,
with
\end{region}

\begin{region}[x=74.7,right=0]{503.1}
\tcode{controller.connect("/dev/ttyUSB0")}
\end{region}

\begin{region}[x=74.7,right=0,pitch=13.2]{521.5}
to open the serial interface. If this fails at startup, it is not tried
again later.
\end{region}

\begin{region}[x=74.7,right=0,pitch=13.53]{561.1}
A simple but practical solution is to extend the program so that the user
can try to activate the game controller at any time during the game by
pressing the \textbf{P} key (for ``peripheral''). This is not yet real hot
plug'n'play, but it is a solution that works in everyday use.
\end{region}

\begin{region}[x=74.7,right=0]{628.9}
\subheading{And this is exactly where you come into play:}
\end{region}

\begin{region}[x=74.7,right=0,pitch=13.2]{655.4}
Adapt your program so that this works. When everything works in the end,
don't forget:
\end{region}

\begin{region}[x=74.7,right=0,pitch=13.2]{695.2}
\code{git commit -a -m "plug and play"}\\
\code{git push}
\end{region}

\nextpage
%% ==================================================================== 10 ==
\begin{region}{69.4}
\heading{Outlook: what awaits you in the next lab handout}
\end{region}

\begin{region}[pitch=13.1]{88.3}
The trickiest part is done -- congratulations! From now on, it is only about
expanding the game further and making it more and more exciting. Here is
what awaits you next time:
\end{region}

\begin{region}[pitch=13.2]{119.5}
\begin{dashlist}
\item There will be bricks that are more robust and have to be hit several
  times.
\item Some bricks have surprising side effects: with some of them you get
  an extra ball, with others additional spaceships suddenly appear or other
  exciting things happen.
\item And the spaceships? Unfortunately, they are no longer as harmless as
  before -- they start dropping bombs. Luckily, there are also bricks that
  give you a protective shield.
\item And \dots\ well, let yourselves be surprised!
\end{dashlist}
\end{region}

\begin{region}[pitch=13.23]{218.7}
\textbf{The nice thing:} You don't have to learn much that is new for this.
Most of it will consist of reorganizing the existing code so that everything
fits together cleanly. And that is completely normal in programming. At the
beginning of a project, you simply cannot foresee every detail -- and that
is exactly what makes it exciting.
\end{region}

\begin{region}{285.6}
\heading{For those who want to experiment on their own}
\end{region}

\begin{region}[pitch=13.17]{304.5}
If you feel like implementing your own features before the next lab
handout, here are a few ideas on how to proceed. The important thing: always
make changes that are as small as possible and that you can test right
away. That way, you reach your goal step by step, without everything
blowing up in your face at the end.
\end{region}

\begin{region}[pitch=13.3]{363.6}
\subheading{Global management of balls and spaceships}
Currently, the lists \code{balls} and \code{ships} are managed locally in
\code{breakout.py}. This means that new balls or spaceships can only be
brought into the game there. A good first exercise would be to move these
lists into the class \code{Game} as static attributes. Then you can access
them from other modules. This is a small change that can be implemented
quickly and tested right away.
\end{region}

\begin{region}[pitch=13.45]{449.9}
\subheading{Bricks with different effects}
In \code{wall.py} there is the matrix \code{Wall.brick}. So far, each entry
has only the value \code{0} (no brick) or \code{1} (normal brick). You could
extend the initialization so that the value \code{2} also occurs there at
random. In the function \code{collision} you can then check what kind of
brick was hit: with \code{1} everything stays as before, with \code{2} you
additionally bring a new ball into the game. Once that works, you will
probably notice quickly: you cannot tell by looking at the bricks what is
special about them. You can change that by adapting the \code{draw}
function.
\end{region}

\begin{region}[pitch=13.55]{550.3}
\subheading{Adding spaceships}
Similarly, you can initialize bricks with the values \code{0}, \code{1},
\code{2}, or \code{3}. A brick with the value \code{3} could, for example,
bring a new spaceship into the game.
\end{region}

\begin{region}[pitch=13.45]{597.1}
\subheading{More robust bricks}
You can extend the range of values further, so that a brick with \code{4}
does not disappear after the first collision but is reset to \code{1}. You
then have to hit it twice to destroy it.
\end{region}

\begin{region}[pitch=13.3]{657.0}
\subheading{Dropping bombs}
Bombs can be implemented much like rockets. You could write a class Bomb
that is structured analogously to the class Rocket. You could also keep a
list of the bombs as a static attribute in the class \code{Game}, so that it
is globally visible.
\end{region}

\begin{region}[pitch=13.1]{723.4}
{\color{aslabblue}
When you try this out: have fun, and remember to test small intermediate
steps regularly! In programming, this is often the decisive trick.
\par}
\end{region}

\begin{region}[pitch=13.35]{755.2}
{\color{aslabblue}
And don't forget to commit your progress to the Git repository. If you
create new files, you have to register them once beforehand with
\code{git add FILENAME} -- after that, they are part of your project just
like the other files.
\par}
\end{region}

\end{document}
